% O ARIMA sendo reestimado: ordens (p, d, q) escolhidas a cada 21 pregões e
% coeficientes reestimados todo pregão (traço da refação de 05/10/2026).
\documentclass[border=6pt]{standalone}
% Preâmbulo comum dos gráficos que explicam os métodos (30/10/2026).
% Dados e números: dados/resultados/metodos_ilustrados_2026-10-30/, gerados por
% ferramentas/MetodosIlustrados.mq5 com as funções do próprio robô. Nenhum
% número é digitado aqui: os do texto vêm de valores.tex (\vX para o texto,
% \pX para coordenada).
\usepackage[T1]{fontenc}
\usepackage{amsmath}
\usepackage{mathptmx}  % Times, como as figuras do TG
\usepackage{pgfplots}
\usepgfplotslibrary{groupplots,dateplot,colormaps,fillbetween}
\usetikzlibrary{arrows.meta,positioning}
\pgfplotsset{compat=1.18}
\definecolor{cmod}{RGB}{0,90,181}     % modelo / ajuste
\definecolor{cdes}{RGB}{220,50,32}    % destaque, limite, valor crítico
\definecolor{cdad}{RGB}{120,120,120}  % dado observado
\newcommand{\arq}[1]{../../dados/resultados/metodos_ilustrados_2026-10-30/#1}
\input{\arq{valores.tex}}
% Figuras 08 a 14 (o robô operando): traço por barra da refação de 05/10/2026,
% resumido por ferramentas/scripts/robo_em_acao.sh. Macros \vR e \pR.
\newcommand{\arqr}[1]{../../dados/resultados/robo_em_acao_2026-10-30/#1}
\input{\arqr{valores.tex}}
\definecolor{cari}{RGB}{230,140,0}   % ARIMA puro
\definecolor{cok}{RGB}{0,140,70}     % compra liberada
\usetikzlibrary{shapes.misc}
\newcommand{\passo}[1]{\tikz[baseline=(n.base)]\node[circle,fill=black!75,text=white,inner sep=1pt,font=\footnotesize\bfseries](n){#1};}
% eixo de datas comum às figuras do traço (jul/2019 a jun/2021)
\pgfplotsset{traco/.style={date coordinates in=x, xmin=2019-07-01, xmax=2021-06-30,
    xtick={2019-07-01,2020-01-01,2020-07-01,2021-01-01,2021-06-30},
    xticklabels={jul/2019,jan/2020,jul/2020,jan/2021,jun/2021}}}
\pgfplotsset{
  base/.style={table/col sep=comma, grid=major, grid style={gray!20},
    tick label style={font=\small}, label style={font=\small},
    title style={font=\small\bfseries, align=center},
    every axis plot/.append style={line width=0.9pt},
    /pgf/number format/use comma, /pgf/number format/1000 sep={.}},
}
% faixa vertical nos pregões do traço em que a condição #1 vale, na cor #2,
% do piso #3 ao teto #4 do eixo
\newcommand{\faixa}[4]{%
  \addplot[name path=topo, draw=none, const plot, forget plot] table[x=data, y expr={(#1) ? #4 : #3}]{\arqr{dia_a_dia.csv}};
  \addplot[name path=piso, draw=none, forget plot] coordinates {(2019-07-01,#3) (2021-06-30,#3)};
  \addplot[#2, forget plot] fill between[of=topo and piso];}

\begin{document}
\begin{tikzpicture}
\begin{groupplot}[base, traco, group style={group size=1 by 3, vertical sep=1.05cm, xticklabels at=edge bottom},
  width=15cm, height=4.3cm]
\nextgroupplot[ylabel={Preço (R\$)}, ymax=32, ymin=10,
  title={\passo{1} O preço da PETR4 na janela (faixas cinza: trechos em que o ADF escolheu $d = 0$)}]
\faixa{\thisrow{d}==0}{black!12}{10}{32}
\addplot[cdad!60!black] table[x=data, y=preco]{\arqr{dia_a_dia.csv}};
\nextgroupplot[ymin=-0.3, ymax=5.8, ytick={0,...,4}, ylabel={Ordem},
  title={\passo{2} As ordens que o AIC escolheu: \vRModelos{} combinações $(p,d,q)$ em dois anos},
  legend style={at={(0.005,0.97)}, anchor=north west, font=\footnotesize, legend columns=2}]
\faixa{\thisrow{d}==0}{black!12}{-0.3}{5.8}
\addplot[cmod, const plot, forget plot] table[x=data, y=p]{\arqr{dia_a_dia.csv}};
\addplot[cari, const plot, dashed, forget plot] table[x=data, y=q]{\arqr{dia_a_dia.csv}};
\addlegendimage{cmod, line width=0.9pt}\addlegendentry{$p$ (retornos passados)}
\addlegendimage{cari, dashed, line width=0.9pt}\addlegendentry{$q$ (erros passados)}
\nextgroupplot[ymin=-0.35, ymax=1.35, ylabel={Coeficiente},
  title={\passo{3} Os coeficientes, reestimados todo pregão por Hannan-Rissanen},
  legend style={at={(0.005,0.97)}, anchor=north west, font=\footnotesize, legend columns=2},
  /pgf/number format/fixed]
\faixa{\thisrow{d}==0}{black!12}{-0.35}{1.35}
\addplot[black!40, forget plot, line width=0.5pt] coordinates {(2019-07-01,0) (2021-06-30,0)};
\addplot[cmod] table[x=data, y=phi1]{\arqr{dia_a_dia.csv}};
\addlegendentry{$\phi_1$}
\addplot[cari] table[x=data, y=theta1]{\arqr{dia_a_dia.csv}};
\addlegendentry{$\theta_1$}
\end{groupplot}
\node[font=\footnotesize, text=black!75, align=left, anchor=north west, text width=14.2cm]
  at ([yshift=-0.75cm]group c1r3.south west) {Como ler: a cada 21 pregões o robô refaz o ADF e a grade do AIC e pode trocar de modelo
  (houve \vRTrocas{} trocas). Entre uma escolha e outra o modelo é o mesmo, mas os coeficientes são reestimados todo dia com a
  janela que andou um pregão. Em \vRDiasPasseio{} dos \vRN{} pregões a escolha foi ARIMA(0,1,0): sem $\phi$ nem $\theta$, a
  previsão é só a constante, o retorno médio da janela. Com $d = 0$ (faixas cinza) o modelo trabalha no log-preço, e o $\phi_1$
  perto de 1 faz quase o papel da diferença.};
\end{tikzpicture}
\end{document}
